% BookID: cd1_b6
% Libro: Cálculo. Primer curso, nivel superior
% Autor: Arizmendi Peimbert, Carrillo Hoyo & Lara Aparicio
% Edición: Addison-Wesley Iberoamericana, 1987

\documentclass[12pt]{article}
\usepackage{ukishima-notas}

\begin{document}

\section*{Arquímedes (287 a 212 a.C.)}

La grandeza del pensamiento científico de Arquímedes ha dominado el panorama de la ciencia durante varios siglos y trasciende hasta nuestros días. Hijo del astrónomo Fidias, nació en Siracusa. Se educó en Alejandría, centro universal del conocimiento. Como si lo anterior fuera poco, se adelantó 2000 años a Newton y Leibniz e inventó el cálculo integral. Encontró que la razón entre la circunferencia y el diámetro de un círculo es igual a $\pi$, cuyo valor estimó entre $3\tfrac{10}{71}$ y $3\tfrac{1}{7}$. Descubrió las leyes de las palancas y creó la ciencia de la hidrostática. Murió a los 75 años cuando los romanos tomaron Siracusa; sus últimas palabras: «¡No moleste a mis círculos!»

\section{Capítulo 1: Números Reales}

\subsection{Introducción}

En el estudio del cálculo tiene importancia primordial el conocimiento de los números reales. Zenón de Elea (siglo~V a.C.) enunció las siguientes paradojas: (1)~\textbf{La dicotomía}: para llegar de $A$ a $B$ hay que pasar primero por la mitad, y así indefinidamente. (2)~\textbf{Aquiles y la tortuga}: Aquiles nunca alcanza a la tortuga pues siempre debe llegar al punto donde ella estaba. (3)~\textbf{La flecha}: en cada instante indivisible la flecha está en reposo, luego no puede moverse. El propósito del capítulo es hacer un modelo numérico de la recta.

\subsection{Números Racionales}

Un \textbf{número racional} es todo número expresable como $p/q$ con $q \neq 0$. Dos expresiones $p/q$ y $r/s$ representan el mismo racional si $ps = qr$. Cuando $p$ y $q$ no tienen factores comunes, $p/q$ es la \textit{mínima expresión}. Las operaciones son:
\[ \frac{p}{q} + \frac{r}{s} = \frac{ps+qr}{qs}, \qquad \frac{p}{q} \cdot \frac{r}{s} = \frac{pr}{qs}. \]

\subsubsection{Orientación de una recta}

Una recta orientada tiene un \textit{sentido positivo} y uno \textit{negativo}. El segmento dirigido de $P$ a $Q$ se denota $\overrightarrow{PQ}$; su longitud se denota $PQ$.

\subsubsection{Identificación de racionales con puntos de la recta}

La $q$-ésima parte de $AB$ se obtiene trazando un rayo desde $A$, marcando $q$ copias de $AU$, uniendo el extremo final $C$ con $B$, y trazando por $U$ una paralela a $CB$; la intersección $U'$ con $AB$ satisface $AU' = AB/q$ ($\triangle AUU' \sim \triangle ABC$). A cada $p/q > 0$ le corresponde el punto a $p$ copias de la $q$-ésima parte de la unidad $\overrightarrow{OP_1}$; si $p/q < 0$, el simétrico respecto al origen.

\subsection{Números Irracionales}

\begin{proposicion}[]{Irracionalidad de $\sqrt{2}$}
No existe ningún número racional cuyo cuadrado sea igual a $2$.
\end{proposicion}

\begin{proof}
Supóngase que $p/q$ en mínima expresión satisface $p^2/q^2 = 2$, es decir, $p^2 = 2q^2$. Entonces $p$ es par: sea $p = 2s$; sustituyendo, $4s^2 = 2q^2$, luego $q^2 = 2s^2$ y $q$ también es par. Contradicción con la mínima expresión.
$\blacksquare$
\end{proof}

\begin{definicion}[]{Número irracional}
Las expresiones decimales no periódicas se denominan \textbf{números irracionales}. Ejemplos: $\sqrt{2}$, $\pi$, $e$, $\sqrt{3}$, $\sqrt{5}$.
\end{definicion}

\subsubsection{Expresiones decimales}

Las expresiones decimales periódicas corresponden a racionales; las no periódicas, a irracionales. Método para convertir periódica a fracción: multiplicar por potencias de $10$ para alinear el período y restar.

\begin{observacion}[]{Convención}
No se consideran expresiones con «colas» de nueves: $2.7\overline{9}$ se identifica con $2.8$.
\end{observacion}

\subsection{La Recta Coordenada}

\begin{definicion}[]{Recta coordenada}
La recta cuyos puntos quedan en correspondencia biunívoca con los números reales (expresiones decimales, periódicas o no) se llama \textbf{recta coordenada}; el punto asociado a $0$ se llama \textbf{origen}.
\end{definicion}

\begin{nota}[]{Sobre $\pm\infty$}
Los símbolos $-\infty$ y $+\infty$ no son números reales; indican la orientación de la recta.
\end{nota}

\subsubsection{Interpretaciones geométricas}

\textbf{Suma:} Los segmentos $\overrightarrow{OA}$ (valor $a$) y $\overrightarrow{AB}$ (valor $b$) componen $\overrightarrow{OB}$, correspondiente a $a+b$.

\textbf{Producto:} Con rectas $L_1, L_2$ que se cortan en $O$ y mismo segmento unitario: unir $U$ (unidad en $L_2$) con $A$ (valor $a$ en $L_1$); trazar por $B$ (valor $b$ en $L_2$) una paralela a $UA$; la intersección $C$ con $L_1$ verifica $OC = ab$ por $\triangle OUA \sim \triangle OBC$.

\textbf{Cociente $a/b$:} Análogamente, unir $A$ con $B$ y trazar por $U$ una paralela; $C$ en $L_1$ verifica $OC = a/b$.

\textbf{Raíz cuadrada:} Con $OA = a$ y $AB = 1$, el círculo de diámetro $OB$ y la perpendicular en $A$ se cortan en $C$ con $OC = \sqrt{a}$ (por $\triangle OAC \sim \triangle ACB$: $x/1 = a/x$).

\subsection{Propiedades de los Números Reales}

\subsubsection{Propiedades de la adición}

\begin{axioma}[A1]{Cerradura}
Para todo $a$ y $b$ reales, $a + b$ es real.
\end{axioma}

\begin{axioma}[A2]{Conmutatividad}
Para todo $a$ y $b$ reales, $a + b = b + a$.
\end{axioma}

\begin{axioma}[A3]{Asociatividad}
Para todo $a$, $b$ y $c$ reales, $(a+b)+c = a+(b+c)$.
\end{axioma}

\begin{axioma}[A4]{Neutro aditivo}
Existe un único $0$ tal que $a + 0 = a$ para todo $a$.
\end{axioma}

\begin{axioma}[A5]{Inverso aditivo}
Para cada $a$ existe un único $-a$ tal que $a + (-a) = 0$.
\end{axioma}

\subsubsection{Propiedades de la multiplicación}

\begin{axioma}[M1]{Cerradura}
Para todo $a$ y $b$ reales, $ab$ es real.
\end{axioma}

\begin{axioma}[M2]{Conmutatividad}
Para todo $a$ y $b$ reales, $ab = ba$.
\end{axioma}

\begin{axioma}[M3]{Asociatividad}
Para todo $a$, $b$ y $c$ reales, $(ab)c = a(bc)$.
\end{axioma}

\begin{axioma}[M4]{Neutro multiplicativo}
Existe un único $1$ tal que $a \cdot 1 = a$ para todo $a$.
\end{axioma}

\begin{axioma}[M5]{Inverso multiplicativo}
Para cada $a \neq 0$ existe un único $a^{-1}$ tal que $a a^{-1} = 1$.
\end{axioma}

\begin{axioma}[D1]{Distributividad}
Para todo $a$, $b$ y $c$ reales, $a(b+c) = ab + ac$.
\end{axioma}

\subsubsection{Orden}

\begin{axioma}[O1]{Tricotomía}
Para cada dos reales $a, b$: exactamente una de $a < b$, $a > b$, $a = b$.
\end{axioma}

\begin{axioma}[O2]{Transitividad}
Si $a < b$ y $b < c$, entonces $a < c$.
\end{axioma}

\begin{axioma}[O3]{Propiedad arquimediana}
Para $a > 0$ y $b$ real, existe $n \in \mathbb{N}$ tal que $na > b$.
\end{axioma}

\begin{axioma}[O4]{}
Si $a, b > 0$, entonces $a + b > 0$.
\end{axioma}

\begin{axioma}[O5]{}
Si $a, b > 0$, entonces $ab > 0$.
\end{axioma}

\begin{axioma}[O6]{Axioma del supremo}
Todo subconjunto no vacío y acotado superiormente de $\mathbb{R}$ tiene supremo.
\end{axioma}

\subsubsection{Propiedades deducidas del orden}

\begin{proposicion}[1]{}
$a > b$ si, y sólo si, $a - b > 0$.
\end{proposicion}

\begin{proof}
$\overrightarrow{OC} = \overrightarrow{OA} + \overrightarrow{BO} = \overrightarrow{BA}$ tiene orientación positiva $\iff C$ es positivo $\iff a - b > 0$.
$\blacksquare$
\end{proof}

\begin{proposicion}[2]{}
Si $a > b$, entonces $a + c > b + c$ para todo $c$.
\end{proposicion}

\begin{proof}
$(a+c) - (b+c) = a - b > 0$.
$\blacksquare$
\end{proof}

\begin{proposicion}[3]{}
Si $a > b$ y $c > d$, entonces $a + c > b + d$.
\end{proposicion}

\begin{proposicion}[4]{}
Si $c > 0$ y $a > b$, entonces $ac > bc$.
\end{proposicion}

\begin{proposicion}[5]{}
$c > 0$ si, y sólo si, $-c < 0$.
\end{proposicion}

\begin{proposicion}[6]{}
Si $c < 0$ y $a > b$, entonces $ac < bc$.
\end{proposicion}

\begin{proposicion}[7]{}
Si $a \neq 0$, entonces $a^2 > 0$.
\end{proposicion}

\begin{proposicion}[8]{}
Si $0 < a < b$ y $0 < c < d$, entonces $ac < bd$.
\end{proposicion}

\begin{proposicion}[9]{}
$1/a$ ($a \neq 0$) tiene el mismo signo que $a$.
\end{proposicion}

\begin{proposicion}[10]{}
Si $a < b$ y ambos tienen el mismo signo, entonces $1/a > 1/b$.
\end{proposicion}

\begin{proof}
$\tfrac{b-a}{ba} = \tfrac{1}{a} - \tfrac{1}{b} > 0$ pues $b - a > 0$ y $ba > 0$ (O5).
$\blacksquare$
\end{proof}

\begin{proposicion}[11]{}
Para $a, b > 0$: $a^2 < b^2$ si, y sólo si, $a < b$.
\end{proposicion}

\begin{nota}[]{}
$x \leq y$ significa $x < y$ o $x = y$. Los signos $<$, $>$ son \textit{estrictos}; $\leq$, $\geq$ son \textit{suaves}.
\end{nota}

\subsection{Intervalos}

\begin{definicion}[]{Intervalo abierto}
Si $a < b$: $(a,b) = \{x \mid a < x < b\}$.
\end{definicion}

\begin{definicion}[]{Intervalo cerrado}
$[a,b] = \{x \mid a \leq x \leq b\}$. Semiabiertos: $(a,b]$ y $[a,b)$. Infinitos: $[a,\infty)$, $(-\infty,a]$, $(a,\infty)$, $(-\infty,a)$.
\end{definicion}

\subsection{Inecuaciones}

Una \textit{solución} de una inecuación es un real $c$ que la satisface al sustituirse. Para inecuaciones de primer grado, se despeja como en ecuaciones invirtiendo la desigualdad al multiplicar/dividir por negativo. Para $ax^2+bx+c < 0$ con raíces $r_1 \leq r_2$: $ax^2+bx+c = a(x-r_1)(x-r_2)$; el signo del producto determina la solución.

\subsection{Valor Absoluto}

\begin{definicion}[]{Valor absoluto}
\[ |a| = \begin{cases} a & \text{si } a \geq 0, \\ -a & \text{si } a < 0. \end{cases} \]
$|a| \geq 0$ y es la longitud del segmento $\overrightarrow{OA}$.
\end{definicion}

\begin{proposicion}[1]{}
$|x|^2 = x^2$.
\end{proposicion}

\begin{proposicion}[2]{}
$|x| = \sqrt{x^2}$ (raíz no negativa).
\end{proposicion}

\begin{proposicion}[3]{}
$|-x| = |x|$.
\end{proposicion}

\begin{proposicion}[4]{}
$|xy| = |x|\,|y|$.
\end{proposicion}

\begin{proposicion}[5]{}
$x \leq |x|$ y $-x \leq |x|$.
\end{proposicion}

\begin{proposicion}[6]{Desigualdad del triángulo}
$|x + y| \leq |x| + |y|$.
\end{proposicion}

\begin{proof}
$|x+y|^2 = x^2+2xy+y^2 \leq |x|^2+2|x||y|+|y|^2 = (|x|+|y|)^2$. Por P.11: $|x+y| \leq |x|+|y|$.
$\blacksquare$
\end{proof}

\begin{proposicion}[7]{}
Para $d > 0$: $|x| < d \iff -d < x < d$.
\end{proposicion}

\begin{proof}
($\Rightarrow$) $x \leq |x| < d$ y $-x \leq |x| < d$, luego $-d < x < d$.
($\Leftarrow$) $(x+d)(d-x) = d^2-x^2 > 0$, luego $x^2 < d^2$, y por P.11 $|x| < d$.
$\blacksquare$
\end{proof}

\begin{observacion}[]{}
La distancia entre $A$ y $B$ es $d(a,b) = |a-b|$; satisface no negatividad ($d(a,b)=0 \iff a=b$), simetría y desigualdad del triángulo $d(a,c) \leq d(a,b)+d(b,c)$.
\end{observacion}

En general: $|x-a| < d \iff x \in (a-d, a+d)$; \quad $|x-a| \leq d \iff x \in [a-d, a+d]$.

\subsection{Elementos Máximo y Mínimo de un Conjunto}

\begin{definicion}[]{Máximo}
$M = \max A$ si $M \in A$ y $M \geq a$ para todo $a \in A$.
\end{definicion}

\begin{definicion}[]{Mínimo}
$m = \min A$ si $m \in A$ y $m \leq a$ para todo $a \in A$.
\end{definicion}

Todo subconjunto no vacío y finito de $\mathbb{R}$ tiene máximo y mínimo. El intervalo $(0,1)$ no tiene ninguno de los dos.

\subsection{El Axioma del Supremo}

\begin{definicion}[1.1]{Supremo}
Sea $A \subseteq \mathbb{R}$ no vacío y acotado superiormente. El número $c = \sup A$ es la \textbf{mínima cota superior} de $A$.
\end{definicion}

\begin{definicion}[1.2]{Ínfimo}
Sea $A \subseteq \mathbb{R}$ no vacío y acotado inferiormente. El número $d = \inf A$ es la \textbf{máxima cota inferior} de $A$.
\end{definicion}

Si $A$ tiene máximo $M$, entonces $M = \sup A$; si tiene mínimo $m$, entonces $m = \inf A$.

\subsubsection{Propiedades características del supremo}

$c = \sup A$ si y sólo si: (1)~$c \geq a$ para todo $a \in A$; y (2)~para todo $d < c$ existe $a \in A$ con $d < a \leq c$.

\section*{Galileo Galilei (1564 a 1642)}

Sentó las bases de la ciencia moderna al relacionar experimentación y teoría. Estudió el movimiento e inició el camino hacia el cálculo. Defendió las ideas de Copérnico; desarrolló el telescopio y descubrió cráteres lunares, manchas solares, satélites de Júpiter y fases de Venus. Propuso la ley de caída libre y sugirió el péndulo en los relojes. Murió el 8 de enero de 1642 en Arcetri bajo arresto domiciliario.

% [Continúa — Capítulo 2: Funciones]

\end{document}
